% =====================================================================
%  VERSI ENGLISH untuk jurnal JETS-ITB (Journal of Engineering and
%  Technological Sciences). Diadaptasi dari nids-01.tex (master, ID).
%  Gaya JETS: APA name-year citation, struktur Introduction / Materials
%  and Methods / Results / Discussion / Conclusion, Calibri-like body.
%  Catatan: template resmi JETS berupa Word (.docx). Berkas .tex ini
%  adalah versi kerja English; kontennya nanti disalin ke template Word
%  dan diberi Style (Heading 1/2/3, Text, Table, Figure, Reference).
% =====================================================================
\documentclass[11pt, a4paper]{article}

% --- Packages ---
\usepackage[utf8]{inputenc}
\usepackage[T1]{fontenc}
\usepackage[english]{babel}
\usepackage[margin=2cm, top=2cm, left=3cm, bottom=2cm, right=1.25cm]{geometry}
\usepackage{amsmath, amsfonts, amssymb}
\usepackage{graphicx}
\usepackage{booktabs}
\usepackage{tabularx}
\usepackage{array}
\usepackage{float}
\usepackage{tikz}
\usepackage{pgfplots}
\pgfplotsset{compat=1.18}
\usetikzlibrary{shapes.geometric, arrows.meta, positioning, fit, backgrounds}
\usepackage{titlesec}
\usepackage{abstract}
\usepackage[round, authoryear]{natbib}   % APA-like name-year citation
\usepackage[hidelinks]{hyperref}

\setlength{\absleftindent}{0pt}
\setlength{\absrightindent}{0pt}
\renewcommand{\abstractnamefont}{\normalfont\bfseries}
\renewcommand{\abstracttextfont}{\normalfont\small}

\titleformat{\section}{\large\bfseries}{\thesection.}{0.5em}{}
\titleformat{\subsection}{\normalfont\bfseries}{\thesubsection}{0.5em}{}
\titleformat{\subsubsection}{\normalfont\itshape}{\thesubsubsection}{0.5em}{}
\setlength{\parindent}{0pt}
\setlength{\parskip}{6pt}

\begin{document}

% =====================================================================
%  FRONT MATTER
% =====================================================================
\begin{center}
    {\LARGE \bfseries Strengthening NIDS Resilience through XGBoost Gain-based Feature Reduction and Adversarial Training against Saliency Map Evasion Attacks \par}
    \vspace{0.4cm}
    {\normalsize Hero Yudo Martono\textsuperscript{1,*}, Iwan Syarif\textsuperscript{1}, Ferry Astika Saputra\textsuperscript{1} \par}
    \vspace{0.2cm}
    {\small \textsuperscript{1}Department of Informatics Engineering, Politeknik Elektronika Negeri Surabaya, Surabaya, Indonesia \par}
    \vspace{0.1cm}
    {\small Email: hero@pens.ac.id; iwanarif@pens.ac.id; ferryas@pens.ac.id \par}
    \vspace{0.1cm}
    {\small \textsuperscript{*}Corresponding author: hero@pens.ac.id \par}
\end{center}

\vspace{0.3cm}

\begin{abstract}
\noindent Machine learning-based Network Intrusion Detection Systems (NIDS) often perform well only under clean laboratory traffic and remain vulnerable to deliberate manipulation during the inference phase. At the same time, feature dimensionality reduction is needed for system efficiency, yet it may enlarge the attack surface. This study evaluates and strengthens the robustness of an XGBoost model whose features have been optimized using the Gain-based Importance method. The CSE-CIC-IDS2018 dataset, comprising 15.7 million network traffic samples across seven attack categories, is used. The proposed methodology has three main stages: (1) feature reduction from 68 to 10 features using the XGBoost Gain metric; (2) simulation of Saliency Map-based evasion attacks with epsilon variations ($\epsilon \in \{0.01, 0.05, 0.1\}$); and (3) implementation of Adversarial Training as a defense mechanism. Because XGBoost is a non-differentiable tree-based model, input sensitivity is approximated numerically through finite differences over the model's probability outputs. Experimental results show that the baseline model suffers a drastic drop in the Matthews Correlation Coefficient (MCC) from 0.9351 to 0.0184 (a 98\% decline) under evasion attack ($\epsilon = 0.1$). The model hardened through Adversarial Training recovers the MCC to 0.9953 on adversarial data while preserving detection integrity on normal traffic (MCC = 0.9347, only 0.04\% below the baseline). A Robustness Ablation Study over four feature configurations (Full/68, Top-15, Top-10, Top-5) confirms that the Top-10 configuration is the optimal sweet spot, balancing computational efficiency and resilience. Furthermore, validation on real network traffic in an AWS cloud environment shows that the robustness pattern remains consistent with the offline results---the baseline degrades under evasion while the robust model holds---demonstrating that the effectiveness of Adversarial Training is not confined to benchmark datasets.

\vspace{0.3cm}
\noindent\textit{\textbf{Keywords:} adversarial training; CSE-CIC-IDS2018; evasion attack; feature reduction; network intrusion detection system; saliency map; XGBoost.}
\end{abstract}

\vspace{0.3cm}

% =====================================================================
%  BODY (diisi bertahap)
% =====================================================================

\section{Introduction}
The rapid growth of network traffic and the increasing sophistication of cyberattacks have made the Network Intrusion Detection System (NIDS) an essential defensive component. The development of machine learning-based NIDS depends heavily on the availability of high-quality public benchmark datasets. Several datasets have become de facto standards for evaluating detection models, ranging from KDD Cup 99 \citep{ref_kdd}, UNSW-NB15 \citep{ref_unsw}, CIC-IDS/CSE-CIC-IDS \citep{ref_cicids}, Bot-IoT \citep{ref_botiot}, and TON\_IoT \citep{ref_toniot}, to CICIoT2023 \citep{ref_ciciot}. Various surveys confirm that machine learning and deep learning methods have become mainstream in modern NIDS research \citep{ref_review_gan, ref_survey_malware, ref_tutorial_gan}.

One classic challenge in NIDS datasets is the extreme class imbalance between normal traffic and attacks. To address this, Generative Adversarial Networks (GAN), introduced by \citet{ref_gan}, are widely used for data synthesis and balancing. Several generative architecture variants have been developed as foundations, including Conditional GAN \citep{ref_cgan}, Wasserstein GAN \citep{ref_wgan}, Deep Convolutional GAN \citep{ref_dcgan}, Self-Attention GAN \citep{ref_sagan}, Auxiliary Classifier GAN \citep{ref_acgan}, and Conditional Tabular GAN \citep{ref_ctgan}.

At the NIDS application level, several studies combine generative models with traditional machine learning classifiers such as Random Forest, Decision Tree, SVM, and XGBoost to improve detection \citep{ref_mari, ref_kumar, ref_ctgan_ids, ref_mcgan, ref_syngan, ref_cegan}. Another group of studies employs deep learning architectures such as CNN, DNN, LSTM, and GRU as downstream classifiers after data generation \citep{ref_tcn, ref_speacgan, ref_qgan, ref_tmggan, ref_park, ref_fewshot}. These approaches have proven capable of improving detection performance under clean laboratory traffic conditions.

However, the same generative capability can be exploited as an offensive weapon. Several studies show that GANs can generate adversarial samples that deceive NIDS models through evasion attacks, on both machine learning-based and deep learning-based classifiers \citep{ref_attackgan, ref_sganids, ref_idsgan}. Such attacks are even effective under low-data regimes \citep{ref_evasiongan}, in machine learning-based evasion schemes \citep{ref_demgan}, in IoT environments \citep{ref_iotevasion}, and in the malware domain \citep{ref_aagan, ref_gleam}. These findings reveal a fundamental security gap: a model that performs excellently on clean data can collapse drastically when its input features are subtly and deliberately manipulated during the inference phase.

To close this gap, defense mechanisms such as Adversarial Training and diffusion-based adversarial purification have been proposed to strengthen the model's decision boundary \citep{ref_diffpurif}. Generative approaches have also evolved toward ensemble/hybrid methods \citep{ref_ewgan, ref_hybridwgan} as well as privacy-preserving and distributed architectures such as differentially private GAN \citep{ref_hegan}, distributed GAN \citep{ref_distgan}, and federated learning \citep{ref_nidsfgpa}. Meanwhile, the demand for computational efficiency drives the adoption of feature dimensionality reduction so that models can be deployed on resource-constrained devices. Ironically, aggressive feature reduction may enlarge the attack surface: by concentrating decisions on a few features, the gradient sensitivity of those features increases, making the model easier to exploit. This tension between efficiency and robustness has not been systematically explored in the literature.

In addition, a review of generative and adversarial NIDS literature reveals two important gaps. \textit{First}, the issue of cross-dataset generalization: recent studies prove that a model excelling on one dataset suffers a drastic performance drop when tested on a dataset from a different network \citep{ref_crossgen, ref_synthtraffic}, indicating that benchmark performance does not necessarily reflect generalization capability. \textit{Second}, and most critically, almost all such studies stop at offline evaluation using static datasets---very few validate model robustness on real network traffic in an actual deployment environment. As a result, claims of robustness against evasion remain largely untested under realistic operational conditions, where differences in feature extraction tools and network characteristics can significantly affect performance. The inherent challenges of training generative models, such as mode collapse and convergence instability, further complicate the reproducibility of results in real environments \citep{ref_challenges}.

Based on these gaps, this study proposes an approach that integrates XGBoost Gain-based feature reduction with Saliency Map-based Adversarial Training, then validates it not only offline on the CSE-CIC-IDS2018 dataset but also on real traffic in an AWS cloud environment. The main contributions of this study are: (1) quantifying the security gap caused by evasion attacks on a feature-reduced tree-based model; (2) demonstrating the effectiveness of Adversarial Training in restoring model robustness without sacrificing accuracy on normal traffic; (3) determining the feature-dimension sweet spot through a Robustness Ablation Study; and (4) validating the consistency of model robustness on real network traffic---an aspect rarely addressed in adversarial NIDS literature.

\section{Materials and Methods}
The methodology of this study is systematically designed to evaluate and improve the robustness of an NIDS against evasion attacks. It comprises six main stages: dataset preprocessing, XGBoost Gain-based feature optimization, adversarial sample generation, adversarial training, robustness ablation study, and model evaluation, as illustrated in Figure~\ref{fig:flowchart}.

\begin{figure}[htbp]
\centering
\includegraphics[width=0.9\linewidth]{nids-figures/fig_flowchart.png}
\caption{research methodology flowchart}
\label{fig:flowchart}
\end{figure}

\subsection{Dataset and Data Preprocessing}
This study uses the CSE-CIC-IDS2018 dataset provided by the Canadian Institute for Cybersecurity (CIC) \citep{ref_cicids}. This dataset was selected because it offers broad attack coverage reflecting modern network scenarios, consisting of ten CSV files with a total size of approximately 6.7 GB. It covers benign traffic profiles and seven main attack categories: Brute-Force (SSH/FTP), DoS, DDoS, Web Attacks, Infiltration, and Botnet.

Given the massive data volume---particularly the 20 February traffic file that reaches 3.8 GB---this study applies a 10\% stratified sampling strategy \citep{ref_stratified} to keep the original attack-label distribution representative in the working subset. The sampling produced 1{,}570{,}577 samples out of a total population of 15{,}705{,}786 samples. Data cleaning was performed systematically through several crucial stages to ensure model integrity:

\begin{enumerate}
    \item \textbf{Handling of Singular Values:} All samples containing missing values (NaN) were permanently removed. Infinity values in numeric features were imputed using the finite maximum of the corresponding column to prevent mathematical failure during model training.
    \item \textbf{Elimination of Biased Features:} Flow-identifier features such as Flow ID, Source/Destination IP, Source/Destination Port, and Protocol were dropped because they do not contribute to attack-pattern generalization and risk introducing training bias. The Timestamp column was also removed to ensure the model learns statistical traffic characteristics rather than temporal ordering.
    \item \textbf{Low-Variance Reduction:} Features with zero variance (constant values) were scanned and discarded because they carry no discriminative information for the classifier.
    \item \textbf{Normalization and Encoding:} Qualitative attack labels were converted to numeric form using \texttt{LabelEncoder}, and numeric features were normalized using \texttt{StandardScaler}/\texttt{MinMax} scaling from scikit-learn \citep{ref_sklearn} to accelerate XGBoost convergence.
\end{enumerate}

Starting from 83 initial technical features, this process yielded a clean data matrix of 68 independent numeric features and one target label, ready for the Gain-based feature selection stage. The distribution of samples before and after preprocessing is summarized in Table~\ref{tab:distribution}.

\begin{table}[H]
\caption{Sample distribution of CSE-CIC-IDS2018 before and after preprocessing}
\label{tab:distribution}
\centering
\small
\begin{tabular}{@{} l r r l @{}}
\toprule
\textbf{Category} & \textbf{Original} & \textbf{10\% Subset} & \textbf{Description} \\ \midrule
Benign        & 13{,}484{,}708 & 1{,}348{,}470 & Normal traffic \\
DDoS / DoS    & 1{,}391{,}082  & 139{,}108     & Volumetric flooding \\
Brute-Force   & 380{,}943      & 38{,}094      & SSH/FTP login \\
Botnet        & 286{,}191      & 28{,}619      & C2 command \\
Web Attacks   & 928           & 93           & SQLi / XSS \\
Infiltration  & 161{,}934      & 16{,}193      & Internal penetration \\ \midrule
\textbf{Total} & \textbf{15{,}705{,}786} & \textbf{1{,}570{,}577} & 83 features (before cleaning) \\ \bottomrule
\end{tabular}
\end{table}

\subsection{XGBoost Gain-based Feature Optimization}
This stage aims to achieve high computational efficiency without sacrificing information critical for anomaly detection. Given the large feature dimensionality (68 columns after preprocessing), Gain-based feature reduction \citep{ref_xgboost} was performed through the following steps: (1) training a baseline XGBoost model on all 68 features with an 80:20 train--test split; (2) extracting the Gain metric, which represents the average contribution of a feature in reducing the loss function at each split of the decision trees; (3) ranking features from highest to lowest Gain; (4) validating via an ablation study by retraining the model on several feature subsets---Full (68), Top-20, Top-15, and Top-10; and (5) selecting the optimal configuration based on a balance of classification performance ($F_1$-score and MCC), model-artifact size, and inference latency. The configuration meeting these three aspects in a balanced manner is designated as the sweet spot.

\subsection{Adversarial Sample Generation (Evasion Attack)}
This stage simulates an attacker attempting to deceive the NIDS through an evasion attack. Adversarial sample generation employs a gradient-based technique \citep{ref_szegedy} to identify model vulnerabilities and apply subtle perturbations to critical features.

\subsubsection{Mathematical Foundation of the Saliency Map}
The Saliency Map approach relies on the notion of output sensitivity to input changes introduced by \citet{ref_saliency}. Let the NIDS be a decision function $f(\mathbf{x})$ with parameters $\theta$, where $\mathbf{x} = [x_1, x_2, \dots, x_d]^T \in \mathbb{R}^d$ is the network-traffic feature vector and $y$ the ground-truth label. The sensitivity of the loss function $L(\theta, \mathbf{x}, y)$ to changes in each feature is measured by the gradient vector:

\begin{equation}
\nabla_{\mathbf{x}} L(\theta, \mathbf{x}, y) = \left[ \frac{\partial L}{\partial x_1}, \frac{\partial L}{\partial x_2}, \dots, \frac{\partial L}{\partial x_d} \right]^T
\end{equation}

The Saliency Map score $S(\mathbf{x}, i)$ for feature $i$ is the absolute value of the partial derivative of the loss with respect to that feature:

\begin{equation}
S(\mathbf{x}, i) = \left| \frac{\partial L(\theta, \mathbf{x}, y)}{\partial x_i} \right|
\end{equation}

A high $S(\mathbf{x}, i)$ indicates that a minimal change to that feature has the most significant effect on the model's decision.

\subsubsection{Gradient Approximation for a Non-Differentiable Model}
It must be emphasized that XGBoost is a tree-based model whose decision boundary is a step function, and is therefore non-differentiable: the partial derivative $\partial L / \partial x_i$ is zero almost everywhere and undefined exactly at the split thresholds. Consequently, the gradient vector in Eq. (1) and Eq. (2) is \textbf{not computed analytically}; instead, it is approximated numerically using the central finite-difference method over the cross-entropy loss constructed from the model's output probabilities (\texttt{predict\_proba} under the \texttt{multi:softprob} objective). This probability function is continuous with respect to input perturbations, so the gradient approximation is well defined. For each feature $i$, the saliency score is computed as:

\begin{equation}
S(\mathbf{x}, i) = \left| \frac{L(\theta, \mathbf{x} + h\mathbf{e}_i, y) - L(\theta, \mathbf{x} - h\mathbf{e}_i, y)}{2h} \right|
\end{equation}

where $\mathbf{e}_i$ is the unit basis vector along dimension $i$, $h = 0.01$ is the perturbation step size, and $L(\cdot)$ is the cross-entropy $L = -\sum_{c} y_c \log p_c(\mathbf{x})$ with $p_c(\mathbf{x})$ the model's probability for class $c$. This approach is conceptually equivalent to a score-based (zeroth-order) attack that requires only access to the model's output probability scores, not its internal gradients. Thus, the vulnerability computation remains mathematically valid even for a non-differentiable tree-based model.

\subsubsection{Epsilon Perturbation Formulation}
Once the most sensitive features are identified, the adversarial sample $\mathbf{x}_{adv}$ is generated by adding a perturbation vector $\boldsymbol{\delta}$ to the original sample $\mathbf{x}$, bounded by the $L_\infty$ norm $\epsilon$:

\begin{equation}
\mathbf{x}_{adv} = \mathbf{x} + \boldsymbol{\delta}, \quad \text{with } \|\boldsymbol{\delta}\|_\infty \le \epsilon
\end{equation}

Following the Fast Gradient Sign Method (FGSM) introduced by \citet{ref_fgsm}, and substituting the approximated saliency for the analytic gradient, the adversarial sample is formed as:

\begin{equation}
\mathbf{x}_{adv} = \mathbf{x} + \epsilon \cdot \text{sign}\left( S(\mathbf{x}) \right)
\end{equation}

On the feature-reduced model, the decision is concentrated on a very limited set of dimensions; the fewer features used, the larger the approximated saliency $S(\mathbf{x}, i)$ on those features tends to be, since the discriminative load concentrates on a narrower space. This mathematical weakness is exploited through the Saliency Map-based attack to produce the adversarial set $\mathbf{x}_{adv}$.

\subsection{Modeling and Adversarial Training}
The NIDS model is built using XGBoost, optimized on the reduced feature subset from Gain-based selection. The objective function is \texttt{multi:softprob} to classify traffic into benign and various attack classes. Key hyperparameters are a maximum tree depth (\texttt{max\_depth}) of 8, a learning rate of 0.1, 200 estimators (\texttt{n\_estimators}), and \texttt{subsample} and \texttt{colsample\_bytree} of 0.8 each to mitigate overfitting.

Adversarial Training is formulated as a Min-Max (saddle-point) optimization following the robustness framework of \citet{ref_madry}. Intuitively, it minimizes the worst-case loss induced by attacker perturbations:

\begin{equation}
\min_{\theta} \mathbb{E}_{(\mathbf{x}, y) \sim \mathcal{D}} \left[ \max_{\|\boldsymbol{\delta}\|_\infty \le \epsilon} L(\theta, \mathbf{x} + \boldsymbol{\delta}, y) \right]
\end{equation}

The practical implementation augments the training set: from the clean training data $\mathcal{D}_{clean}$, an adversarial subset $\mathcal{D}_{adv}$ is generated using the Saliency Map method, and a new training set $\mathcal{D}_{robust}$ is formed by combining clean and adversarial samples at an 80:20 ratio:

\begin{equation}
\mathcal{D}_{robust} = \mathcal{D}_{clean} \cup \mathcal{D}_{adv}
\end{equation}

Retraining XGBoost on $\mathcal{D}_{robust}$ adjusts the split thresholds of the decision trees so that regions around the original samples---previously vulnerable to an $\epsilon$-shift---are now covered within the correct decision region, improving robustness without sacrificing accuracy on normal traffic. The complete adversarial training workflow is illustrated in Figure~\ref{fig:advtrain}.

\begin{figure}[htbp]
\centering
\includegraphics[width=0.95\linewidth]{nids-figures/fig_advtrain.png}
\caption{systematic adversarial training workflow}
\label{fig:advtrain}
\end{figure}

\subsection{Robustness Ablation Study}
The Robustness Ablation Study evaluates the trade-off between feature-dimension efficiency and the model's security against evasion attacks. Testing isolates the feature variable into four hierarchical levels based on the XGBoost Gain ranking, as detailed in Table~\ref{tab:ablation_config}.

\begin{table}[H]
\caption{Feature configurations for the Robustness Ablation Study}
\label{tab:ablation_config}
\centering
\small
\begin{tabular}{@{}c c p{0.5\linewidth}@{}}
\toprule
\textbf{Configuration} & \textbf{Count} & \textbf{Description} \\ \midrule
$C_1$ (Full)   & 68 & Full-dimension comparison baseline \\
$C_2$ (Top-15) & 15 & Moderately informative feature subset \\
$C_3$ (Top-10) & 10 & Sweet-spot hypothesis (compact \& precise) \\
$C_4$ (Top-5)  & 5  & Extreme compression (maximum efficiency) \\ \bottomrule
\end{tabular}
\end{table}

Each configuration $C_k$ ($k \in \{1,2,3,4\}$) is trained with two approaches---baseline (no defense) and robust (with Adversarial Training)---then tested with evasion samples $\mathbf{x}_{adv}$ parameterized by $\epsilon$. The goal is to identify the sweet spot $\hat{C}$ that maximizes a combined utility function:

\begin{equation}
\hat{C} = \arg\max_{C_k} \left( \text{MCC}(C_k, \mathbf{x}_{adv}) - \lambda \cdot |C_k| \right)
\end{equation}

where $\text{MCC}(C_k, \mathbf{x}_{adv})$ measures detection robustness on manipulated traffic, $|C_k|$ is the number of features, and $\lambda$ is a dimension-efficiency penalty weight.

\subsection{Model Evaluation and Performance Comparison}
Evaluation uses a $2\times 2$ cross-testing scenario comparing the standard (baseline) model against the hardened (robust) model on both clean and manipulated traffic, as summarized in Table~\ref{tab:scenario_2x2}.

\begin{table}[H]
\caption{Evaluation scenario matrix for the NIDS model ($2\times 2$)}
\label{tab:scenario_2x2}
\centering
\small
\begin{tabular}{@{} l l l l p{0.36\linewidth} @{}}
\toprule
\textbf{Scenario} & \textbf{Model} & \textbf{Train} & \textbf{Test} & \textbf{Evaluation objective} \\ \midrule
S1 & Baseline & Clean & Clean       & Initial performance under normal conditions \\
S2 & Baseline & Clean & Manipulated & Measure evasion vulnerability (security gap) \\
S3 & Robust   & Robust & Clean      & Verify accuracy is preserved on normal traffic \\
S4 & Robust   & Robust & Manipulated & Prove improved resilience against attacks \\ \bottomrule
\end{tabular}
\end{table}

Because CSE-CIC-IDS2018 exhibits extreme class imbalance, the Matthews Correlation Coefficient (MCC) \citep{ref_matthews} is chosen as the primary metric. Unlike conventional accuracy, which can be misleading on imbalanced data, MCC accounts for all four quadrants of the confusion matrix proportionally, making it more reliable for imbalanced classification \citep{ref_chicco}:

\begin{equation}
\text{MCC} = \frac{(TP \times TN) - (FP \times FN)}{\sqrt{(TP + FP)(TP + FN)(TN + FP)(TN + FN)}}
\end{equation}

The MCC ranges over $[-1, +1]$, where $+1$ denotes perfect prediction, $0$ random guessing, and $-1$ total disagreement. As supporting metrics, Precision and $F_1$-score are also reported:

\begin{equation}
\text{Precision} = \frac{TP}{TP + FP}, \qquad
F_1 = 2 \times \frac{\text{Precision} \times \text{Recall}}{\text{Precision} + \text{Recall}}
\end{equation}

where $\text{Recall} = TP/(TP + FN)$. A performance drop in S2 reveals the security gap, whereas recovery of MCC and $F_1$-score in S4 demonstrates the success of model hardening through Adversarial Training.

\section{Results}

\subsection{Preprocessing Results and Subset Description}
The preprocessing transformed the massive raw CSE-CIC-IDS2018 dataset ($\approx 6.7$ GB) into a clean, proportionally balanced subset. The 10\% stratified sampling preserved the original attack-category ratios. During data cleaning, 2{,}861 samples containing undefined values (NaN/Missing) were eliminated, and Infinity values in quantitative variables such as \texttt{Flow Bytes/s} were imputed with the finite maximum of the respective column. Removing six flow-identifier features (Flow ID, Source/Destination IP, Ports, Protocol) and the Timestamp variable prevented overfitting to specific network topology, while a variance scan discarded eight constant (zero-variance) features. From 83 initial technical features, preprocessing produced a clean data matrix of 68 independent numeric features plus one target label.

\subsection{XGBoost Gain Ranking and Sensitive-Feature Exploration}
Training the baseline model on all 68 features showed that the Gain contribution is not evenly distributed but concentrated on a few dominant attributes. Table~\ref{tab:feature_importance} lists the top-20 features by Gain score. Cumulatively, the top-10 features account for more than 85\% of the total accumulated Gain, empirically confirming that most of the 68 features are redundant or contribute minimal discriminative information.

\begin{table}[H]
\caption{Top-20 features ranked by XGBoost Gain-based importance}
\label{tab:feature_importance}
\centering
\small
\begin{tabular}{@{}cl cl@{}}
\toprule
\textbf{Rank} & \textbf{Feature} & \textbf{Rank} & \textbf{Feature} \\ \midrule
1  & Fwd Seg Size Min   & 11 & Bwd Pkt Len Std \\
2  & URG Flag Cnt       & 12 & Bwd Seg Size Avg \\
3  & Tot Bwd Pkts       & 13 & Fwd IAT Std \\
4  & Fwd Act Data Pkts  & 14 & Bwd IAT Max \\
5  & Fwd Pkt Len Max    & 15 & Bwd Pkt Len Max \\
6  & Fwd Pkt Len Mean   & 16 & Fwd Header Len \\
7  & Bwd Pkt Len Mean   & 17 & Fwd Pkts/s \\
8  & Init Bwd Win Byts  & 18 & Fwd URG Flags \\
9  & TotLen Bwd Pkts    & 19 & Fwd IAT Tot \\
10 & Init Fwd Win Byts  & 20 & Fwd Pkt Len Std \\ \bottomrule
\end{tabular}
\end{table}

A gradient-based analysis further distinguished global importance (XGBoost Gain) from local sensitivity (approximated Saliency Map). The evaluation showed that when the model is forced to decide from only 10 features, the saliency magnitude on those features increases sharply compared with the 68-feature model. The feature \textit{Init Fwd Win Byts} recorded the highest saliency score, making it the most sensitive to perturbation---even though it ranks only 10th in the Gain metric. Technically, \textit{Init Fwd Win Byts} represents the initial forward TCP window size, which an attacker can manipulate in the TCP header without breaking packet functionality, making it a realistic evasion target.

\subsection{Feature-Selection Ablation Results}
Following the ablation procedure, the XGBoost model was retrained on four feature configurations. Table~\ref{tab:ablation_selection} compares classification performance, model-artifact size, and inference speed.

\begin{table}[H]
\caption{Feature-selection ablation: performance and efficiency comparison}
\label{tab:ablation_selection}
\centering
\small
\begin{tabular}{@{} l c c c c c @{}}
\toprule
\textbf{Configuration} & \textbf{Features} & \textbf{$F_1$-score} & \textbf{ROC-AUC} & \textbf{Model Size} & \textbf{Inference (10k)} \\ \midrule
Full Features & 68 & 0.9798 & 0.9911 & 7.17 MB & 0.2400 s \\
Top-20        & 20 & 0.9795 & 0.9900 & 6.82 MB & 0.2266 s \\
Top-15        & 15 & 0.9717 & 0.9883 & 5.91 MB & 0.1930 s \\
\textbf{Top-10} & \textbf{10} & \textbf{0.9717} & \textbf{0.9865} & \textbf{5.65 MB} & \textbf{0.1826 s} \\ \bottomrule
\end{tabular}
\end{table}

The $F_1$-score decreases only 0.83\% from Full (0.9798) to Top-10 (0.9717), while the model size shrinks 21.2\% (7.17~MB to 5.65~MB) and inference time drops 23.9\% (0.2400~s to 0.1826~s per 10{,}000 samples). Because Top-10 matches the Top-15 $F_1$-score with a more compact and faster model, it was selected as the sweet spot.

\subsection{Comparative Evaluation of the Four Scenarios (S1--S4)}
Table~\ref{tab:results_2x2} presents the quantitative results of the four cross-testing scenarios on the Top-10 subset with perturbation $\epsilon = 0.10$.

\begin{table}[H]
\caption{Quantitative evaluation results for scenarios S1, S2, S3, and S4}
\label{tab:results_2x2}
\centering
\small
\begin{tabular}{@{} l l l c c c @{}}
\toprule
\textbf{Scenario} & \textbf{Model} & \textbf{Test data} & \textbf{MCC} & \textbf{$F_1$-score} & \textbf{Precision} \\ \midrule
S1 (Baseline)      & Baseline & Clean       & 0.9351 & 0.9729 & 0.9720 \\
S2 (Vulnerability) & Baseline & Manipulated & 0.0184 & 0.7539 & 0.6905 \\
S3 (Integrity)     & Robust   & Clean       & 0.9347 & 0.9727 & 0.9718 \\
S4 (Robustness)    & Robust   & Manipulated & \textbf{0.9953} & \textbf{0.9985} & \textbf{0.9985} \\ \bottomrule
\end{tabular}
\end{table}

Under evasion attack, the baseline MCC collapses from 0.9351 (S1) to 0.0184 (S2), a 98\% decline that constitutes the security gap. After Adversarial Training, the robust model recovers to MCC = 0.9953 on manipulated data (S4), an increase of $+0.9769$ over S2. Meanwhile, the robust model on clean data (S3) retains MCC = 0.9347, only 0.04\% below the baseline S1, indicating that hardening does not degrade normal-traffic detection. Notably, the $F_1$-score in S2 remains at 0.7539 despite the near-zero MCC, because the weighted $F_1$ is still inflated by correct classification of the majority (benign) class, whereas MCC more sensitively captures the collapse on the minority (attack) classes.

The per-class behavior across the four scenarios is visualized in the confusion matrices of Figure~\ref{fig:confusion}. In S2, nearly all attack classes are misclassified as benign (the collapse), whereas S4 restores a clean diagonal, confirming the recovery. A grouped-bar summary of the four metrics is shown in Figure~\ref{fig:groupedbar}.

\begin{figure}[htbp]
\centering
\includegraphics[width=\linewidth]{nids-figures/evaluation_confusion_2x2_en.png}
\caption{confusion matrices for the 2x2 evaluation scenarios (XGBoost Top-10, $\epsilon=0.1$)}
\label{fig:confusion}
\end{figure}

\begin{figure}[htbp]
\centering
\includegraphics[width=0.85\linewidth]{nids-figures/evaluation_grouped_bar_2x2_en.png}
\caption{comprehensive metric comparison across scenarios S1--S4}
\label{fig:groupedbar}
\end{figure}

\subsection{Robustness Ablation Results: Feature-Dimension Sweet Spot}
The Robustness Ablation Study evaluated the effect of feature pruning on resilience under evasion ($\epsilon = 0.10$). Table~\ref{tab:robustness_ablation} reports the adversarial-data MCC for the baseline (S2) and robust (S4) models across the four configurations, together with the recovery achieved by Adversarial Training.

\begin{table}[H]
\caption{Robustness ablation: adversarial-data MCC before (S2) and after (S4) Adversarial Training per configuration ($\epsilon = 0.10$)}
\label{tab:robustness_ablation}
\centering
\small
\begin{tabular}{@{} l c c c c @{}}
\toprule
\textbf{Configuration} & \textbf{Features} & \textbf{MCC$_{base+adv}$ (S2)} & \textbf{MCC$_{rob+adv}$ (S4)} & \textbf{Recovery} \\ \midrule
$C_1$ (Full)   & 68 & $-0.047$ & \textbf{0.997} & $+1.044$ \\
$C_2$ (Top-15) & 15 & $-0.006$ & 0.992 & $+0.998$ \\
$C_3$ (Top-10) & 10 & \textbf{0.018} & \textbf{0.995} & $+0.977$ \\
$C_4$ (Top-5)  & 5  & $-0.070$ & 0.873 & $+0.943$ \\ \bottomrule
\end{tabular}
\end{table}

All baseline models collapse under attack (MCC near zero or negative), whereas every robust model recovers to MCC $> 0.87$. The Top-5 configuration ($C_4$) shows the lowest robustness (MCC S4 = 0.873), indicating that five features are insufficient to form a stable decision boundary. Configurations $C_1$, $C_2$, and $C_3$ differ only marginally in MCC S4 (0.997, 0.992, and 0.995 respectively; $< 0.5\%$), while model size grows from 5.65~MB to 7.17~MB. The Top-10 configuration ($C_3$) is the only one that keeps the baseline adversarial MCC positive (0.018) while achieving MCC S4 = 0.995, confirming it as the optimal sweet spot.

\subsection{Real-Traffic Validation in the AWS Cloud Environment}
To verify that the offline findings are not merely laboratory artifacts, the baseline and robust models were tested on real network traffic generated in an Amazon Web Services (AWS) cloud environment. The testing infrastructure was built inside a Virtual Private Cloud (VPC, $10.3.0.0/16$) comprising three EC2 nodes: an Attacker node (public subnet) running Hydra, Slowloris, and \texttt{nping}; a Target node (private subnet) running SSH and HTTP (Nginx) services; and an Analyzer node performing feature extraction and inference. Traffic to the Target was captured using \texttt{tcpdump}, stored to Amazon S3, and extracted into flow features using the NFStream library. Because NFStream does not natively expose two of the Top-10 TCP window features (\textit{Init Fwd Win Byts} and \textit{Init Bwd Win Byts}), a custom NFStream plugin was developed to parse the TCP window value from the first packet of each flow direction, so that all 10 features could be extracted intact. The testing topology is shown in Figure~\ref{fig:aws}.

\begin{figure}[htbp]
\centering
\includegraphics[width=\linewidth]{nids-figures/fig_aws.png}
\caption{real-traffic testing infrastructure on AWS: the attacker node (public subnet) launches attacks against the target node (private subnet); traffic is captured at the target and routed through Amazon S3 to the analyzer node for feature extraction and inference}
\label{fig:aws}
\end{figure}

The test replicated the $2\times 2$ matrix (RT-S1 to RT-S4) equivalent to the offline scenarios. Table~\ref{tab:real_traffic} compares real-traffic results with offline results on the Top-10 configuration.

\begin{table}[H]
\caption{Offline vs. real-traffic (AWS) performance comparison on the Top-10 configuration}
\label{tab:real_traffic}
\centering
\small
\begin{tabular}{@{} l c c c c @{}}
\toprule
\textbf{Scenario} & \textbf{MCC (Offline)} & \textbf{MCC (Real)} & \textbf{$F_1$ (Offline)} & \textbf{$F_1$ (Real)} \\ \midrule
RT-S1 (Base + Clean)     & 0.9351 & 0.164  & 0.9729 & 0.855 \\
RT-S2 (Base + Evasion)   & 0.0184 & $-0.070$ & 0.7539 & 0.633 \\
RT-S3 (Robust + Clean)   & 0.9347 & 0.339  & 0.9727 & 0.581 \\
RT-S4 (Robust + Evasion) & 0.9953 & 0.389  & 0.9985 & 0.626 \\ \bottomrule
\end{tabular}
\end{table}

Although the absolute MCC values decrease on real traffic, the relative behavioral pattern is preserved. The baseline degrades when exposed to evasion (RT-S1 $\rightarrow$ RT-S2: MCC $0.164 \rightarrow -0.070$, falling below random guessing), whereas the robust model instead improves (RT-S3 $\rightarrow$ RT-S4: MCC $0.339 \rightarrow 0.389$) and records perfect Precision ($1.0$). Diagnostic analysis of the feature distributions attributes the absolute drop to feature mismatch between extractors: \textit{Fwd Seg Size Min} in NFStream ($\approx 66$ bytes, including header) differs from the CICFlowMeter definition at training time ($\approx 18$ bytes), a deviation of $z \approx +6.2$; and the TCP window value on modern EC2 instances ($\approx 62$ KB) is far larger than in the CIC-IDS2018 recording ($\approx 8$ KB), a deviation of $z \approx +3.3$ on \textit{Init Fwd/Bwd Win Byts}.

\section{Discussion}

\subsection{Security Gap and the Effectiveness of Adversarial Training}
The experimental results reveal a fundamental vulnerability of the feature-reduced baseline model. Although the baseline achieves a high MCC of 0.9351 under normal conditions (S1), its performance collapses to 0.0184 under a subtle perturbation of $\epsilon = 0.10$ (S2)---a 98\% decline. This confirms that concentrating the decision on a small set of features, while efficient, sharpens the model's local sensitivity and enlarges the attack surface. The disparity between the near-zero MCC and the still-moderate $F_1$-score (0.7539) in S2 also demonstrates why MCC is the more trustworthy metric on imbalanced data: the weighted $F_1$ remains inflated by the majority (benign) class even as detection of the attack classes fails almost entirely.

Adversarial Training addresses this weakness effectively. On clean traffic, the robust model (S3, MCC = 0.9347) is virtually identical to the baseline (S1), only 0.04\% lower---evidence that augmenting the training set with adversarial samples does not induce a ``paranoid'' model or harm normal-traffic detection. On manipulated traffic, the robust model recovers to MCC = 0.9953 (S4), isolating the attacker's gradient-shift within the redrawn decision boundary.

\subsection{Interpreting the S4 $>$ S1 Phenomenon: Transfer vs. White-Box Attack}
The observation that MCC in S4 (0.9953) exceeds S1 (0.9351) warrants careful interpretation. It arises because the adversarial samples were generated from the saliency of the \textit{baseline} model, not the robust model. This constitutes a \textbf{transfer attack} (grey-box with respect to the robust model): perturbations designed to fool the baseline become ineffective against a robust model trained to withstand exactly those patterns, while clean data retains richer natural variation than the stereotypical $\epsilon \cdot \text{sign}(S(\mathbf{x}))$ perturbation. It must be acknowledged honestly that, because of this transfer setup, the S4 robustness value is potentially an \textbf{optimistic estimate} of the model's true robustness. Under a stricter \textbf{white-box adaptive} scenario---where the attacker generates perturbations directly from the robust model's own saliency---the observed robustness would likely be lower. This study did not perform such white-box adaptive evaluation; the limitation is stated explicitly and set as a primary direction for future work. Accordingly, the robustness claim here is valid within the scope of defense against transfer/static evasion attacks, not adaptive white-box attacks.

\subsection{Feature-Dimension Sweet Spot and Practical Implications}
The Robustness Ablation Study confirms Top-10 ($C_3$) as the optimal balance: it reduces model size by 21.2\% (7.17~MB to 5.65~MB) while maintaining robustness (MCC S4 = 0.995) practically identical to the full model (0.997) and slightly above Top-15 (0.992). Extreme pruning to five features ($C_4$) degrades robustness to 0.873, indicating that five dimensions cannot form a stable decision boundary. These findings carry practical weight for modern network-security architectures. First, the compact $\approx 5.65$~MB model can be embedded directly into resource-constrained edge devices such as routers and IoT gateways, removing the dependency on centralized cloud processing. Second, in serverless architectures (e.g., AWS Lambda) where cost scales with execution time and memory, using only 10 features keeps inference very fast and memory-light, reducing operational cost at enterprise scale. Third, the robust model's resilience (MCC = 0.9953) mitigates false-alarm fatigue, ensuring the NIDS is not easily paralyzed by deliberately crafted TCP/IP-variable shifts.

The holistic radar comparison across five metrics (MCC, $F_1$-score, Precision, Recall, Accuracy), shown in Figure~\ref{fig:radar}, reinforces this reading: the near-full, symmetric polygons of S1, S3, and S4 indicate balanced high performance, whereas the sharply shrunken S2 polygon---collapsing on the MCC axis yet still wide on Accuracy and Recall---visually confirms the misleading nature of conventional metrics on imbalanced data. The similar polygon shapes of S1, S3, and S4 further emphasize that Adversarial Training restores robustness without distorting the model's detection profile under normal conditions.

\begin{figure}[htbp]
\centering
\includegraphics[width=0.7\linewidth]{nids-figures/evaluation_radar_2x2_en.png}
\caption{multi-metric radar comparison of scenarios S1--S4}
\label{fig:radar}
\end{figure}

\subsection{Generalization to Real Traffic and Mitigation Strategies}
The real-traffic validation delivers a nuanced but important message. While the absolute MCC drops on AWS traffic, the relative pattern is preserved---the baseline degrades under evasion (MCC $0.164 \rightarrow -0.070$) whereas the robust model improves (MCC $0.339 \rightarrow 0.389$) with perfect Precision. This confirms that the core contribution, the effectiveness of Adversarial Training, is robust across test domains, not merely a benchmark artifact. At the same time, it exposes a generalization challenge rarely discussed in the literature: strong benchmark performance does not guarantee real-world generalization, chiefly because of differences in feature-extraction tooling and network environment. The drop is not a fundamental model failure but a manifestation of distribution shift, driven by the extractor-definition mismatch ($z \approx +6.2$ on \textit{Fwd Seg Size Min}) and the environmental difference in TCP window size ($z \approx +3.3$ on \textit{Init Fwd Win Byts}).

To close this gap and make the system reliably deployable, three concrete mitigation strategies are recommended. \textit{First}, \textbf{extractor calibration}: mathematically aligning the production extractor (NFStream) with the training extractor (CICFlowMeter)---for example, correcting \textit{Fwd Seg Size Min} from the NFStream basis ($\approx 66$ bytes, including header) to the CICFlowMeter definition ($\approx 18$ bytes) via an affine (offset/scaling) transform, which addresses the largest deviation without retraining. \textit{Second}, \textbf{domain adaptation / transfer learning}: fine-tuning the model's decision layer using a small amount of labeled real traffic from the target environment to align the decision boundary with the operational feature distribution, particularly for sensitive features such as \textit{Init Fwd Win Byts}. \textit{Third}, \textbf{multi-extractor training}: incorporating extraction variability from multiple libraries (CICFlowMeter, NFStream, and Zeek) into the training set so that the model becomes invariant to extractor-implementation differences. These strategies are complementary---calibration offers a quick fix without retraining, domain adaptation aligns the model to a specific environment, and multi-extractor training provides structural cross-environment robustness.

\subsection{Comparison with Prior Work}
Unlike the mainstream of adversarial NIDS research, which combines generative models with classifiers and stops at offline evaluation \citep{ref_mari, ref_kumar, ref_park, ref_tmggan}, this study emphasizes the tension between feature efficiency and adversarial robustness on a tree-based model, and validates it on real cloud traffic. Whereas prior evasion studies largely demonstrate attack effectiveness \citep{ref_attackgan, ref_idsgan, ref_demgan}, the present work quantifies the security gap and then closes it through Adversarial Training \citep{ref_madry}, additionally reporting the honest generalization gap on real traffic---an aspect consistent with cross-dataset concerns raised by \citet{ref_crossgen} but rarely examined under live-traffic conditions.

\section{Conclusion}
This study developed and evaluated an NIDS optimization approach that combines XGBoost Gain-based feature reduction with Saliency Map gradient-based Adversarial Training. Four main conclusions follow. First, the feature-reduced baseline model is highly vulnerable to evasion attacks: although it attains MCC = 0.9351 under normal conditions (S1), performance collapses to MCC = 0.0184 under a subtle perturbation ($\epsilon = 0.10$, S2)---a 98\% decline. Second, Adversarial Training based on Min-Max optimization restores detection reliability to MCC = 0.9953 (S4, an increase of $+0.9769$ over S2) while preserving integrity on normal traffic (S3, MCC = 0.9347; only 0.04\% below the baseline). Third, through the Robustness Ablation Study, the Top-10 configuration ($C_3$) proved to be the optimal sweet spot, cutting model size to 5.65~MB (a 21.2\% reduction) while maintaining high resilience (MCC S4 = 0.995); the Top-5 configuration degraded to MCC S4 = 0.873, confirming that five features are insufficient for a stable decision boundary. Fourth, real-traffic validation in the AWS cloud showed that the robustness pattern remains consistent with the offline results---the baseline degrades under evasion while the robust model holds---even though the absolute MCC decreases due to feature mismatch between the NFStream and CICFlowMeter extractors and differences in the TCP window environment. This consistency confirms that the effectiveness of Adversarial Training is robust across test domains, while also revealing a generalization challenge seldom discussed in the literature.

To broaden the scope and implications of this study, several directions are recommended for future research. First, \textbf{white-box adaptive attack evaluation}: the current robustness assessment is limited to a transfer-attack scheme (perturbations generated from the baseline model), so the robust model's robustness may be optimistic; future work should generate adversarial samples directly from the robust model's own saliency to obtain a worst-case estimate. Second, \textbf{black-box attack testing} (e.g., Boundary Attack or query-based attacks) where the attacker has no direct access to the target model's gradients. Third, \textbf{extending classifier types}: applying and comparing this gradient-based hardening strategy on other tree-based architectures (LightGBM, CatBoost) and compressed deep learning models (quantized neural networks). Fourth, \textbf{in-the-wild testbed}: deploying the 5.65~MB model on real edge hardware (e.g., Raspberry Pi or a smart router) to measure real-time packet-processing throughput and power consumption. Finally, implementing and quantitatively evaluating the three real-traffic mitigation strategies (extractor calibration, domain adaptation, and multi-extractor training) to close the benchmark-to-deployment generalization gap.

\section*{Acknowledgement}
The authors thank Politeknik Elektronika Negeri Surabaya (PENS) for the research support and facilities provided during this study.

\section*{Conflicts of Interest}
The authors declare no conflict of interest.

% =====================================================================
%  REFERENCES (APA name-year, alphabetical) — Task #7
% =====================================================================
\renewcommand{\refname}{References}
\begin{thebibliography}{99}
\setlength{\itemsep}{2pt}

\bibitem[Alabsi et al., 2023]{ref_ctgan_ids}
Alabsi, B. A., Anbar, M., \& Rihan, S. D. A. (2023). Conditional tabular generative adversarial based intrusion detection system for detecting DDoS and DoS attacks on the Internet of Things networks. \textit{Sensors}, 23(12), 5644.

\bibitem[Aldhaheri \& Alhuzali, 2023]{ref_sganids}
Aldhaheri, S., \& Alhuzali, A. (2023). SGAN-IDS: Self-attention-based generative adversarial network against intrusion detection systems. \textit{Sensors}, 23(18), 7796.

\bibitem[Al-Ajlan \& Ykhlef, 2024]{ref_review_gan}
Al-Ajlan, M., \& Ykhlef, M. (2024). A review of generative adversarial networks for intrusion detection systems: Advances, challenges, and future directions. \textit{Computers, Materials \& Continua}, 81(2), 2053--2076.

\bibitem[Alshehri et al., 2025]{ref_hybridwgan}
Alshehri, M. S., Saidani, O., Malwi, W. A., Asiri, F., Latif, S., Khattak, A. A., \& Ahmad, J. (2025). A hybrid Wasserstein GAN and autoencoder model for robust intrusion detection in IoT. \textit{Computer Modeling in Engineering \& Sciences}, 143(3), 3899--3920.

\bibitem[Alsaedi et al., 2020]{ref_toniot}
Alsaedi, A., Moustafa, N., Tari, Z., Mahmood, A. N., \& Anwar, A. (2020). TON\_IoT telemetry dataset: A new generation dataset of IoT and IIoT for data-driven intrusion detection systems. \textit{IEEE Access}, 8, 165130--165150.

\bibitem[Aceto et al., 2024]{ref_synthtraffic}
Aceto, G., Giampaolo, F., Guida, C., Izzo, S., Pescap\`{e}, A., Piccialli, F., \& Prezioso, E. (2024). Synthetic and privacy-preserving traffic trace generation using generative AI models for training network intrusion detection systems. \textit{Journal of Network and Computer Applications}, 229, 103926.

\bibitem[Arjovsky et al., 2017]{ref_wgan}
Arjovsky, M., Chintala, S., \& Bottou, L. (2017). Wasserstein GAN. \textit{arXiv preprint} arXiv:1701.07875.

\bibitem[Babu \& Rao, 2023]{ref_mcgan}
Babu, K. S., \& Rao, Y. N. (2023). MCGAN: Modified conditional generative adversarial network for class imbalance problems in network intrusion detection system. \textit{Applied Sciences}, 13(4).

\bibitem[Cantone et al., 2024]{ref_crossgen}
Cantone, M., Marrocco, C., \& Bria, A. (2024). Machine learning in network intrusion detection: A cross-dataset generalization study. \textit{IEEE Access}, 12, 144491--144509.

\bibitem[Chen \& Guestrin, 2016]{ref_xgboost}
Chen, T., \& Guestrin, C. (2016). XGBoost: A scalable tree boosting system. In \textit{Proceedings of the 22nd ACM SIGKDD International Conference on Knowledge Discovery and Data Mining} (pp. 785--794).

\bibitem[Chicco \& Jurman, 2020]{ref_chicco}
Chicco, D., \& Jurman, G. (2020). The advantages of the Matthews correlation coefficient (MCC) over F1 score and accuracy in binary classification evaluation. \textit{BMC Genomics}, 21(1), 6.

\bibitem[de Araujo-Filho et al., 2023]{ref_tcn}
de Araujo-Filho, P. F., Naili, M., Kaddoum, G., Fapi, E. T., \& Zhu, Z. (2023). Unsupervised GAN-based intrusion detection system using temporal convolutional networks and self-attention. \textit{IEEE Transactions on Network and Service Management}, 20(4), 4951--4963.

\bibitem[Devadiga et al., 2023]{ref_gleam}
Devadiga, D., Jin, G., Potdar, B., Koo, H., Han, A., Shringi, A., Singh, A., Chaudhari, K., \& Kumar, S. (2023). GLEAM: GAN and LLM for evasive adversarial malware. In \textit{2023 14th International Conference on Information and Communication Technology Convergence (ICTC)} (pp. 53--58).

\bibitem[Ding et al., 2024]{ref_tmggan}
Ding, H., Sun, Y., Huang, N., Shen, Z., \& Cui, X. (2024). TMG-GAN: Generative adversarial networks-based imbalanced learning for network intrusion detection. \textit{IEEE Transactions on Information Forensics and Security}, 19, 1156--1167.

\bibitem[Dunmore et al., 2023]{ref_survey_malware}
Dunmore, A., Jang-Jaccard, J., Sabrina, F., \& Kwak, J. (2023). Generative adversarial networks for malware detection: A survey. \textit{arXiv preprint} arXiv:2302.08558.

\bibitem[Educative, 2025]{ref_challenges}
Educative. (2025). \textit{What are the challenges in training GAN?} Retrieved September 10, 2025, from https://www.educative.io/answers/what-are-the-challenges-in-training-gan

\bibitem[Goodfellow et al., 2014]{ref_gan}
Goodfellow, I., Pouget-Abadie, J., Mirza, M., Xu, B., Warde-Farley, D., Ozair, S., Courville, A., \& Bengio, Y. (2014). Generative adversarial nets. In \textit{Advances in Neural Information Processing Systems} (pp. 2672--2680).

\bibitem[Goodfellow et al., 2015]{ref_fgsm}
Goodfellow, I. J., Shlens, J., \& Szegedy, C. (2015). Explaining and harnessing adversarial examples. In \textit{International Conference on Learning Representations (ICLR)}.

\bibitem[Hassan et al., 2023]{ref_hegan}
Hassan, U., Chen, D., Cheung, S.-C., \& Chuah, C.-N. (2023). HE-GAN: Differentially private GAN using Hamiltonian Monte Carlo based exponential mechanism. In \textit{Proceedings of IEEE ICASSP} (pp. 3186--3190).

\bibitem[Huang et al., 2022]{ref_tutorial_gan}
Huang, Y., Fields, K. G., \& Ma, Y. (2022). A tutorial on generative adversarial networks with application to classification of imbalanced data. \textit{Statistical Analysis and Data Mining: The ASA Data Science Journal}, 15(5), 543--552.

\bibitem[Jiang et al., 2023]{ref_distgan}
Jiang, X., Zhang, Y., Zhou, X., \& Grossklags, J. (2023). Distributed GAN-based privacy-preserving publication of vertically-partitioned data. \textit{Proceedings on Privacy Enhancing Technologies}, 2023(2), 50--69.

\bibitem[Kohavi, 1995]{ref_stratified}
Kohavi, R. (1995). A study of cross-validation and bootstrap for accuracy estimation and model selection. In \textit{Proceedings of the 14th International Joint Conference on Artificial Intelligence (IJCAI)} (pp. 1137--1143).

\bibitem[Koroniotis et al., 2019]{ref_botiot}
Koroniotis, N., Moustafa, N., Sitnikova, E., \& Turnbull, B. (2019). Towards the development of realistic botnet dataset for IoT networks. \textit{IEEE Access}, 7, 94529--94541.

\bibitem[Kumar \& Sinha, 2023]{ref_kumar}
Kumar, V., \& Sinha, D. (2023). Synthetic attack data generation model applying generative adversarial network for intrusion detection. \textit{Computers \& Security}, 125, 103054.

\bibitem[Madry et al., 2018]{ref_madry}
Madry, A., Makelov, A., Schmidt, L., Tsipras, D., \& Vladu, A. (2018). Towards deep learning models resistant to adversarial attacks. In \textit{International Conference on Learning Representations (ICLR)}.

\bibitem[Mari et al., 2023]{ref_mari}
Mari, A.-G., Zinca, D., \& Dobrota, V. (2023). Development of a machine-learning intrusion detection system and testing of its performance using a generative adversarial network. \textit{Sensors}, 23(3), 1315.

\bibitem[Matthews, 1975]{ref_matthews}
Matthews, B. W. (1975). Comparison of the predicted and observed secondary structure of T4 phage lysozyme. \textit{Biochimica et Biophysica Acta (BBA) - Protein Structure}, 405(2), 442--451.

\bibitem[Mbow et al., 2024]{ref_iotevasion}
Mbow, M., Roman, R., Takahashi, T., \& Sakurai, K. (2024). Evading IoT intrusion detection systems with GAN. In \textit{2024 19th Asia Joint Conference on Information Security (AsiaJCIS)} (pp. 48--55).

\bibitem[Merzouk et al., 2025]{ref_diffpurif}
Merzouk, M. A., Beurier, E., Yaich, R., Boulahia-Cuppens, N., Cuppens, F., \& Khomh, F. (2025). Diffusion-based adversarial purification for intrusion detection. In \textit{IFIP Annual Conference on Data and Applications Security and Privacy} (pp. 351--370). Springer.

\bibitem[Mirza \& Osindero, 2014]{ref_cgan}
Mirza, M., \& Osindero, S. (2014). Conditional generative adversarial nets. \textit{arXiv preprint} arXiv:1411.1784.

\bibitem[Moustafa \& Slay, 2015]{ref_unsw}
Moustafa, N., \& Slay, J. (2015). UNSW-NB15: A comprehensive data set for network intrusion detection systems. In \textit{MILCOM} (pp. 1--6).

\bibitem[Neto et al., 2023]{ref_ciciot}
Neto, E. C. P., Dadkhah, S., Ferreira, R., Zohourian, A., Lu, R., \& Ghorbani, A. A. (2023). CICIoT2023: A real-time dataset and benchmark for large-scale attacks in IoT environment. \textit{Sensors}, 23(13).

\bibitem[Odena et al., 2017]{ref_acgan}
Odena, A., Olah, C., \& Shlens, J. (2017). Conditional image synthesis with auxiliary classifier GANs. In \textit{Proceedings of the 34th International Conference on Machine Learning} (PMLR, Vol. 70, pp. 2642--2651).

\bibitem[Park et al., 2023]{ref_park}
Park, C., Lee, J., Kim, Y., Park, J.-G., Kim, H., \& Hong, D. (2023). An enhanced AI-based network intrusion detection system using generative adversarial networks. \textit{IEEE Internet of Things Journal}, 10(3), 2330--2345.

\bibitem[Pedregosa et al., 2011]{ref_sklearn}
Pedregosa, F., Varoquaux, G., Gramfort, A., Michel, V., Thirion, B., Grisel, O., Blondel, M., Prettenhofer, P., Weiss, R., Dubourg, V., Vanderplas, J., Passos, A., Cournapeau, D., Brucher, M., Perrot, M., \& Duchesnay, E. (2011). Scikit-learn: Machine learning in Python. \textit{Journal of Machine Learning Research}, 12, 2825--2830.

\bibitem[Radford et al., 2015]{ref_dcgan}
Radford, A., Metz, L., \& Chintala, S. (2015). Unsupervised representation learning with deep convolutional generative adversarial networks. \textit{arXiv preprint} arXiv:1511.06434.

\bibitem[Rahman et al., 2024]{ref_syngan}
Rahman, S., Pal, S., Mittal, S., Chawla, T., \& Karmakar, C. (2024). Syn-GAN: A robust intrusion detection system using GAN-based synthetic data for IoT security. \textit{Internet of Things}, 26, 101212.

\bibitem[Rahman et al., 2023]{ref_qgan}
Rahman, M. A., Shahriar, H., Clincy, V., Hossain, M. F., \& Rahman, M. (2023). A quantum generative adversarial network-based intrusion detection system. In \textit{2023 IEEE 47th Annual Computers, Software, and Applications Conference (COMPSAC)} (pp. 1810--1815).

\bibitem[Randhawa et al., 2023]{ref_evasiongan}
Randhawa, R. H., Aslam, N., Alauthman, M., \& Rafiq, H. (2023). Evasion generative adversarial network for low data regimes. \textit{IEEE Transactions on Artificial Intelligence}, 4(5), 1076--1088.

\bibitem[Riaz et al., 2025]{ref_ewgan}
Riaz, R., Han, G., Shaukat, K., Khan, N. U., Zhu, H., \& Wang, L. (2025). A novel ensemble Wasserstein GAN framework for effective anomaly detection in industrial Internet of Things environments. \textit{Scientific Reports}, 15(1), 26786.

\bibitem[Sharafaldin et al., 2018]{ref_cicids}
Sharafaldin, I., Lashkari, A. H., \& Ghorbani, A. A. (2018). Toward generating a new intrusion detection dataset and intrusion traffic characterization. In \textit{ICISSP} (pp. 108--116).

\bibitem[Simonyan et al., 2014]{ref_saliency}
Simonyan, K., Vedaldi, A., \& Zisserman, A. (2014). Deep inside convolutional networks: Visualising image classification models and saliency maps. In \textit{Workshop at International Conference on Learning Representations (ICLR)}.

\bibitem[Szegedy et al., 2014]{ref_szegedy}
Szegedy, C., Zaremba, W., Sutskever, I., Bruna, J., Erhan, D., Goodfellow, I., \& Fergus, R. (2014). Intriguing properties of neural networks. In \textit{International Conference on Learning Representations (ICLR)}.

\bibitem[Tavallaee et al., 2009]{ref_kdd}
Tavallaee, M., Bagheri, E., Lu, W., \& Ghorbani, A. A. (2009). A detailed analysis of the KDD Cup 99 data set. In \textit{Proceedings of CISDA}.

\bibitem[Trung et al., 2024]{ref_aagan}
Trung, D. M., Khoa, N. H., Duy, P. T., Pham, V.-H., \& Cam, N. T. (2024). AAGAN: Android malware generation system based on generative adversarial network. \textit{International Journal of Semantic Computing}, 14(1).

\bibitem[Wang et al., 2024a]{ref_idsgan}
Wang, D., Wang, X., \& Fei, J. (2024). IDS-GAN: Adversarial attack against intrusion detection based on generative adversarial networks. In \textit{2024 5th International Conference on Computer Vision, Image and Deep Learning (CVIDL)} (pp. 1130--1134).

\bibitem[Wang et al., 2024b]{ref_nidsfgpa}
Wang, J., Yang, K., \& Li, M. (2024). NIDS-FGPA: A federated learning network intrusion detection algorithm based on secure aggregation of gradient similarity models. \textit{PLoS ONE}, 19(10), e0308639.

\bibitem[Xu et al., 2019]{ref_ctgan}
Xu, L., Skoularidou, M., Cuesta-Infante, A., \& Veeramachaneni, K. (2019). Modeling tabular data using conditional GAN. In \textit{Advances in Neural Information Processing Systems} (Vol. 32). Curran Associates, Inc.

\bibitem[Xu et al., 2025a]{ref_demgan}
Xu, D., Lv, Y., Wang, M., Zheng, B., Zhao, J., \& Yu, J. (2025). DEMGAN: A machine learning-based intrusion detection system evasion scheme. \textit{Computers, Materials \& Continua}, 84(1), 1731--1746.

\bibitem[Xu et al., 2025b]{ref_fewshot}
Xu, C., Zhan, Y., Chen, G., Wang, Z., Liu, S., \& Hu, W. (2025). Elevated few-shot network intrusion detection via self-attention mechanisms and iterative refinement. \textit{PLoS ONE}, 20(1), e0317713.

\bibitem[Yang et al., 2023]{ref_speacgan}
Yang, H., Xu, J., Xiao, Y., \& Hu, L. (2023). SPE-ACGAN: A resampling approach for class imbalance problem in network intrusion detection systems. \textit{Electronics}, 12(15), 3323.

\bibitem[Yang et al., 2025]{ref_cegan}
Yang, Y., Liu, X., Wang, D., Sui, Q., Yang, C., Li, H., Li, Y., \& Luan, T. (2025). A CE-GAN based approach to address data imbalance in network intrusion detection systems. \textit{Scientific Reports}, 15(1), 7916.

\bibitem[Zhang et al., 2019]{ref_sagan}
Zhang, H., Goodfellow, I., Metaxas, D., \& Odena, A. (2019). Self-attention generative adversarial networks. In \textit{Proceedings of the 36th International Conference on Machine Learning} (PMLR, Vol. 97, pp. 7354--7363).

\bibitem[Zhao et al., 2021]{ref_attackgan}
Zhao, S., Li, J., Wang, J., Zhang, Z., Zhu, L., \& Zhang, Y. (2021). AttackGAN: Adversarial attack against black-box IDS using generative adversarial networks. \textit{Procedia Computer Science}, 187, 128--133.

\end{thebibliography}

\end{document}
